\section{Classical and quantum realizations from raw kernels}\label{sec:realizations}
The two realizations below are verified after the general theorem. Recurrent witness preparation, when present, belongs to the charged productive primitive; it is not a reset made freely available to other models. Neither supplies a premise used to prove it. Their conditional physical evolution between calls may be noncontracting, and their acquired laws may be singular. Both tasks score the post-command audit. A target which must remain unchanged is a different experiment; the nonintervening directional result is retained in the companion.

\subsection*{A common paid physical interface}
The gate port $G$ costs one per call and has success probability at least $p_0>0$, conditional on every physical and visible past. A success enables exactly one productive call at port $A$, also of cost one. Commitment is forbidden at $G$ and allowed after a productive report. Rejection at $A$ returns the protocol to $G$; it need not reset the hidden state. The two ports have disjoint physical command domains and no port tag is a free intercall input to an autonomous observer. Suppose the actual lower witness has gate success probability at most $p_1\le1$. Then
\begin{equation}\label{eq:gate-constants}
 \kappa_-=1+1/p_1,\qquad \kappa_+=1+1/p_0,
 \qquad B_0=B=2.
\end{equation}
Indeed every productive call consumes a gate success. A terminating path has no unconsumed success. Predictable counting gives $\E_*A_\tau\le p_1\E_*G_\tau$, so $\E_*C_\tau\ge\kappa_-\E_*A_\tau$. The probability of $j$ failures after any gate start is at most $(1-p_0)^j$, giving conditional expected gate count at most $1/p_0$ and whole-attempt cost at most $\kappa_+$. Both mandatory ports are visited before commitment, and Lemma~\ref{lem:roles} gives two live roles. Consequently the construction requires at least two event labels or three state labels when it has one output. If the witness has constant success $p_0=p_1$, then $\E_*C_\tau=\kappa\E_*A_\tau$ with $\kappa=1+1/p_0$. In that matched case budget feasibility is exactly $N\ge\kappa$ with the respective state thresholds. Even when $p_0=1$, the cost is two, not one.

\begin{theorem}[A relatively calibrated classical realization]\label{thm:classical}
Let a hidden symbol $E$ take values in a finite set of size $q$. A finite group $\mathcal G$ acts transitively on that set. The allowed productive commands apply $g\in\mathcal G$ and then, within the same call, a fixed positive normalized joint report--audit kernel $J_{gE}(dy,dz)$ with bounded audit in $H$. There is no unknown intervening operation between that command and its audit instrument. Evolution between port interactions is otherwise arbitrary and causal.

The genuinely prepared comparison has uniform input and joint law $K_0=q^{-1}\sum_eJ_e$. Its full report is its mark. Suppose the operating family contains a single witness whose conditional input weights, after every reached pre-command history, satisfy
\[
 (1-\eta_r)/q\le\Prb_*(E=e\mid h)\le(1+\eta_r)/q,
 \qquad 0\le\eta_r<1.
\]
Suppose one parameter-independent implementation chooses its command independently of the hidden current symbol conditional on its visible past, with conditional weights
\[
 (1-\eta_u)/|\mathcal G|\le\lambda_g\le(1+\eta_u)/|\mathcal G|,
 \qquad0\le\eta_u<1,
\]
uniformly over all models and surviving histories. Use the gate interface above. Then Theorem~\ref{thm:transfer} holds with identity exports, zero additive defects, and
\begin{equation}\label{eq:classical-relative}
 l_-=1-\eta_r,\quad u_-=1+\eta_r,
 \qquad l_+=1-\eta_u,\quad u_+=1+\eta_u.
\end{equation}
The comparison need not be an operating member. Fixed nonzero $\eta_r,\eta_u<1$ preserve the all-scale resource curve of Corollary~\ref{cor:uniform} whenever its acquired geometry holds.
\end{theorem}
\begin{proof}
A chosen permutation preserves the two componentwise bounds on the witness weights. Multiplying by each positive $J_e$ and summing gives the full command-conditioned joint-law order
$(1-\eta_r)K_0\le P_{h,g}\le(1+\eta_r)K_0$.
Take $P^0_{h,g}=K_0$ and the identity mark. Conditional moments given its full report are exactly \eqref{eq:derived}, so all lower certificates follow from Proposition~\ref{prop:matrix-order}. This holds for every chosen command and every competing policy, not merely for a randomizing policy.

For any fixed incoming symbol, uniform averaging over a transitive group gives the uniform distribution of its image: each image has the same number of preimages in the group. Averaging also over any incoming distribution preserves that fact. The actual componentwise command bounds and positivity therefore put the full joint upper-program law between $1-\eta_u$ and $1+\eta_u$ times $K_0$. Independence of the command from the hidden current input is needed for this conditional averaging; it is stated as a physical implementation property, not inferred from marginal coin frequencies. The upper accepts on the resulting report, not on a discarded group choice. Exports are the identity because the audit is the post-command audit. Group selection is within-call work, and only the gate or productive routing label or final decoder value persists. The gate proof establishes the same-witness cost lower, common cost upper and both live roles. The general theorem now applies.
\end{proof}

For a concrete acquired geometry take $q=2$, with identity and sign reversal. Independently draw $(Y,T)$ from a fixed law on $K\times[0,1]$, report these coordinates, and audit
\begin{equation}\label{eq:classical-mark}
 Z=Y+TEe_1
\end{equation}
in a sufficiently large compact convex carrier. The uniform preparation gives $f=Y,w=T^2$. Choose a $d$-regular location measure and $\Prb(T\le s)=s^\alpha$ to verify \eqref{eq:rectangle} directly; choose a finite mixture of such laws to verify Theorem~\ref{thm:stratified}; add a physically prepared point mass at $(x_*,0)$ to obtain Proposition~\ref{prop:atom}. A biased fresh witness and imperfect command probabilities satisfy \eqref{eq:classical-relative} with fixed nonzero constants. On $T=0$ the task is exactly deterministic even under these errors. On $T>0$ the two-point audit moment matrix has rank two, while at $T=0$ it has rank one. This is an actual changing-rank task measure, not a formal limiting parameter set.

The command group and its bit-level sampler are part of the program description and workspace account. A high-cardinality group does not provide free randomization or a free persistent seed. In the binary realization a single fresh sign coin suffices within a productive call. Conditional posterior continuity of the classical raw kernels is governed by Proposition~\ref{prop:continuity}, including likelihood zeros.

\begin{theorem}[A relatively calibrated noncommuting quantum realization]\label{thm:quantum}
Let the productive physical system have finite dimension $d\ge2$. Its available commands include the $d^2$ Weyl conjugations, followed within the same productive call by a fixed positive normalized quantum report--audit instrument $\mathcal J$, with classical joint law $\mathcal J(\rho)(dy,dz)$ and $z\in H$. The physical task is again the post-command audit. Put $\omega=I/d$ and $K_0=\mathcal J(\omega)$.

Suppose one actual witness obeys, conditionally before every selected productive command,
\[
 (1-\eta_r)\omega\preccurlyeq\rho\preccurlyeq(1+\eta_r)\omega,
 \qquad0\le\eta_r<1.
\]
Suppose the common implementation uses current Weyl weights between $(1-\eta_u)/d^2$ and $(1+\eta_u)/d^2$, independently of the current hidden input conditional on the past, with $0\le\eta_u<1$. With the same gate certificates, all conclusions of Theorem~\ref{thm:classical} hold with the constants \eqref{eq:classical-relative}. The operating witness need not be exactly maximally mixed. Incoming states for the upper may be pure, and post-report ranks may change on positive acquired mass. No common eigenbasis or normalized-filter contraction is assumed.
\end{theorem}
\begin{proof}
Let $W_{p,q}$ be the shift by $q$ followed by the phase $\exp(2\pi ipj/d)$ on the $j$th basis vector. On a matrix unit $|j\rangle\langle k|$, averaging the phase kills off-diagonal terms; averaging shifts makes the diagonal uniform. Hence
\[
 d^{-2}\sum_{p,q}W_{p,q}\rho W_{p,q}^*=I/d
\]
for every density matrix. This elementary randomization identity is standard \cite{AMTW}; we use its complete positive-instrument consequence, not a new state-randomization claim.

Each chosen conjugation fixes $\omega$ and preserves the witness Loewner order. Positivity of the full report--audit instrument therefore gives lower joint-law order relative to $K_0$ with factors $1\pm\eta_r$. Positivity means that the image of every positive operator is a positive classical measure, including both report and future audit. It consequently compares the full observable report and task, not just the report marginal.

For the upper, the componentwise weight bounds give
\[
 (1-\eta_u)\omega
 \preccurlyeq\sum_{p,q}\lambda_{p,q}W_{p,q}\rho W_{p,q}^*
 \preccurlyeq(1+\eta_u)\omega
\]
for every incoming state, including pure states. Apply the same positive instrument. Proposition~\ref{prop:matrix-order} supplies the task certificates, and identity export, counted live roles and the gate costs finish the verification. The lower witness has a relative input bound around $\omega$; we do not describe that hypothesis as absence of an input-rank condition. The absence of a floor concerns the post-report states and task strata, where the following example has genuine rank changes.
\end{proof}

Here is a raw qubit instrument giving the same acquired geometry as \eqref{eq:classical-mark} but noncommuting predictive states. Draw $(y,t)$ with any stated law $\nu_0$ on $K\times[0,1]$, choose an equatorial unit vector $n=(\cos\phi,\sin\phi,0)$ and $0\le\gamma(y,t)\le1$, and set
\[
 E_v=(I+v\gamma n\cdot\sigma)/2,\qquad L_v=\sqrt{E_v},\qquad v=\pm1.
\]
The report is $(y,t,v)$ and the unnormalized post-report state is $L_v\rho L_v$. At commitment measure $\sigma_z$, obtain $e=\pm1$, and score $Z=y+te e_1$. Under $\omega=I/2$, each $v$ has mass $1/2$ and the normalized posterior is $E_v$. Its $z$ component is zero, so the future audit sign is fair. Thus the reference report law is $\nu_0$ times a fair $v$, and $f=y,w=t^2$, derived directly from the instrument.

The eigenvalues of $E_v$ are $(1\pm\gamma)/2$. Taking $\gamma=1$ on a positive-mass set and $\gamma<1$ on another gives actual pure and full-rank posterior strata. Varying $\phi$ produces noncommuting effects. For a continuous compact realization one may take $\gamma(t)=1$ for $t\le1/2$ and $\gamma(t)=3/2-t$ for $t\ge1/2$, and a continuous nonconstant angle. The two-sided task order remains uniform across these strata; Proposition~\ref{prop:continuity} supplies actual-mass-weighted predictive-update and report continuity. A point mass at $t=0$ also gives a changing-rank task matrix and an exact zero-loss acquired component.

The qubit is part of the physical experiment, not a quantum memory register owned by the finite classical observer. The pulse and instrument form one productive call in this convention; their duration, finite sampling precision and arithmetic are separately charged. An interface requiring a separate pulse call must add its cost and any intervening routing state. Both realizations retain a strictly positive paid gate floor $p_0$.

\subsection*{Predictive realization rather than an assumed parameter manifold}
For the finite classical kernels and finite-dimensional quantum instruments whose entries depend continuously on compact report coordinates, the primitive continuity criterion in Proposition~\ref{prop:raw-quotient} supplies compact positive-evidence history patches. Their executable predictive quotient is therefore Borel by Theorem~\ref{thm:quotient}. Fixed continuous causal evolutions between calls are included and need not contract. For a completely arbitrary merely Borel hidden evolution, the risk theorem still uses the stated raw certificates, but we do not assert the compact-exhaustion quotient theorem without its regularity hypotheses. The unknown model's family of posteriors is never stored by the common upper program.
